\section{Naravna dedukcija in Curry Howardova korespondenca}\label{natded}

% V tem poglavju bomo najprej predstavili pravila sklepanja, na katerih temelji naravna dedukcija oziroma ``standardna'' logika in s tem bomo tudi predstavili notacijo, ki bo uporabljena v tem diplomskem delu.

V tem poglavju bomo predstavili pravila izvajanja programov, ki so v korespondenci z pravili sklepanja pri naravni dedukcji. S tem bomo tudi predstavili notacijo, ki jo bomo uporabljali v tem diplomskem delu.





\pravilo{Produktni tip} 
Pravilo vpeljave za konjunkcijo pravi, če sta izjavi \(A\) in \(B\) resnični, potem je tudi izjava \(A\) in \(B\) resnična. To lahko zapišemo kot
%
\begin{prooftree}
    \AXC{\(A\)}
    \AXC{\(B\)}
    \RL{\(\land I\).}
    \BIC{\(A \land B\)}
\end{prooftree}

To pravilo je v korespondenci z pravilom vpeljave za produktni tip:
%
\begin{prooftree}
    \AXC{\(N: A\)}
    \AXC{\(M: B\)}
    \RL{\(\times I\).}
    \BIC{\((N, M) : A \times B\)}
\end{prooftree}

Pravila sklepanja imajo obliko 
%
\begin{prooftree}
    \AXC{\(\mathcal{H}_1\)}
    \AXC{\(\mathcal{H}_2\)}
    \AXC{\(\ldots\)}
    \AXC{\(\mathcal{H}_n\)}
    \RL{ime,}
    \QuaternaryInfC{\(\mathcal{S}\)}
\end{prooftree}
%
kjer sodbam \(\mathcal{H}_1, \ldots \mathcal{H}_n\) pravimo premise, pod črto pa sobi \(\mathcal{S}\) pravimo \(sklep\)~\cite[Poglavje 3.1]{Jazbec2022Primerjava}. Desno od črte zapiešmo ime pravila, ki je navadno okrajšano.

Del sodbe je tudi kontekst, ki ga navadno označujemo z \(\G\) ali \(\D\). Kontekst je nabor elementov z njihovimi tipi \(x_1 : A_1, \ldots x_n : A_n\)~\cite[Section 1.1]{rijke2022introduction}.
Sodba je torej oblike \(\G \vdash N: A\), kjer znaku \(vdash\) pravimo \emph{rampa} (ang.~turnstile). 
Pravilo vpeljave za produktni tip je tako oblike
%
\begin{prooftree}
    \AXC{\(\G \vdash N: A\)}
    \AXC{\(\G \vdash M: B\)}
    \RL{\(\times I\).}
    \BIC{\(\G \vdash (N, M) : A \times B\)}
\end{prooftree}
%
To preberemo kot: Če iz konteksta \(\G\) lahko dobimo element \(N\) tipa \(A\) in če iz konteksta \(\G\) lahko dobimo \(M\) tipa \(B\), potem lahko iz konteksta \(\G\) dobimo par \((N, M)\) tipa \(A \times B\).

Pomembno pravilo pri delu s kontekstom je pravilo ošibitve (ang.~weakening rule), ki pravi, da če iz konteksta \(\G\) lahko dobimo \(x: A\), potem lahko to dobimo tudi v razširjenem kontekstu \(\G, y: B\), se pravi v kontekstu, ki mu dodamo nek element.
\begin{prooftree}
    \AXC{\(\G \vdash x: A\)}
    \RL{\(W\)}
    \UIC{\(\G, y: B \vdash x: A\)}
\end{prooftree}

Imamo dve pravili uporabe za produktni tip --- levo in desno.
\begin{equation*}
    \begin{bprooftree}
        \AXC{\(\G \vdash N : A \times B\)}
        \RL{\(\times E_L\)}
        \UIC{\(\G \vdash \fst~N : A\)}
    \end{bprooftree}
    \qquad
    \begin{bprooftree}
        \AXC{\(\G \vdash N : A \times B\)}
        \RL{\(\times E_D\)}
        \UIC{\(\G \vdash \snd~N : A\)}
    \end{bprooftree}
\end{equation*}

Poleg pravil vpeljav iz uporabe imamo tudi pravila računanja (ang.~computation rule), ki v primeru produktnega tipa pove, kako se obnašata funkciji \(\fst\) in \(\snd\)
\begin{align*}
    \fst (N, M) \longrightarrow N  \\
    \snd (N, M) \longrightarrow M 
\end{align*}




\pravilo{Pravilo aksioma}
Preprosto pravilo je pravilo aksioma, ki pravi, da če imamo v konteksu vrednost \(x\) tipa \(A\) jo lahko uporabimo.
%
\begin{prooftree}
    \AXC{}
    \RL{\emph{Ax}}
    \UIC{\(\G, x : A \vdash x : A\)}
\end{prooftree}
%
Za to ne potrebujemo nobenih premis, torej je zgornja črta prazna.

\pravilo{Funkcijski tip}
Malce bolj zahtevno pravilo vpeljave ima funkcijski tip.
%

\begin{prooftree}
    \AXC{\(\G, x : A \vdash M : B\)}
    \RL{\(\rightarrow I\)}
    \UIC{\(\G \vdash \la x. M : A \rightarrow B\)}
\end{prooftree}
%
Če imamo na voljo posamezen element \(x\) tipa \(A\) in lahko dobimo \(M\), potem lahko sestavimo funkcijo, ki elementu \(x  : A\) priredi \(M : B\). Tu je \(M\) lahko odvisen od \(x\), saj pride iz konteksta, ki vsebuje \(x : A\).
Oznaka \(\la x. M\) pove, da je \(x\) vezana spremenljivka v izrazu \(M\) in predstavlja funkcijo, ki elementu \(x\) priredi vrednost \(M\). Pogosto tudi eksplicitno povemo tip vezane spremenlivke, torej pišemo \(\la x : A. M: B\) namesto \(\la x.M : A \rightarrow B\).

Pravilo uporabe za funkcijski tip je
\begin{prooftree}
    \AXC{\(\G \vdash M : A \rightarrow B\)}
    \AXC{\(\G \vdash N : A \)}
    \RL{\( \rightarrow E\).}
    \BIC{\(\G \vdash M~N : B\)}
\end{prooftree}

Pravilo računanja za funkcijski tip je
\begin{equation*}
    (\la x. M) N \longrightarrow M[N/x].
\end{equation*}
%
Oznaka \(M[N/x]\) pove, da zamenjamo vse pojavitve \(x\) v \(M\) z \(N\).


\begin{primer}
    Poskusimo s temi pravili sestaviti element tipa \(A \rightarrow B \rightarrow (A \times B)\).
    \begin{prooftree}
        \AXC{?}
        \UIC{\(\vdash ? : A \rightarrow B \rightarrow (A \times B)\)}
    \end{prooftree}
    Ker želimo dobiti funkcijo, bomo morali uporabiti pravilo vpeljave za funkcijo.
    \begin{prooftree}
        \AXC{?}
        \UIC{\(x: A \vdash ? : B \rightarrow A \times B\)}
        \RL{\(\rightarrow I\)}
        \UIC{\(\vdash \la x. ? : A \rightarrow B \rightarrow (A \times B)\)}
    \end{prooftree}
    To nam razkrije že nekaj več kako bo izgledal končni term, namreč vemo, da bo \(x\) vezana spremenljivka.
    \begin{prooftree}
        \AXC{?}
        \UIC{\(x: A, y: B \vdash ? : A \times B\)}
        \RL{\(\rightarrow I\)}
        \UIC{\(x: A \vdash \la y. ? : B \rightarrow A \times B\)}
        \RL{\(\rightarrow I\)}
        \UIC{\(\vdash \la x. \la y. ? : A \rightarrow B \rightarrow (A \times B)\)}
    \end{prooftree}
    Sedaj želimo dobiti nekaj tipa \(A \times B\). Edina možnost, ki jo imamo je, da uporabimo pravilo vpeljave za produktni tip.
    \begin{prooftree}
        \AXC{?}
        \UIC{\(x:A, y: B \vdash x: A\)}
        \AXC{?}
        \UIC{\(x:A, y: B \vdash y: B\)}
        \RL{\(\times I\)}
        \BIC{\(x: A, y: B \vdash (x, y) : A \times B\)}
        \RL{\(\rightarrow I\)}
        \UIC{\(x: A \vdash \la y. (x, y): B \rightarrow A \times B\)}
        \RL{\(\rightarrow I\)}
        \UIC{\(\vdash \la x. \la y. (x, y) : A \rightarrow B \rightarrow (A \times B)\)}
    \end{prooftree}
    %
    Na vrhu uporabimo pravilo aksioma in s tem zaključimo izpeljavo.
    \begin{prooftree}
        \AXC{}
        \RL{\(Ax\)}
        \UIC{\(x:A, y: B \vdash x: A\)}
        \AXC{}
        \RL{\(Ax\)}
        \UIC{\(x:A, y: B \vdash y: B\)}
        \RL{\(\times I\)}
        \BIC{\(x: A, y: B \vdash (x, y) : A \times B\)}
        \RL{\(\rightarrow I\)}
        \UIC{\(x: A \vdash \la y. (x, y): B \rightarrow A \times B\)}
        \RL{\(\rightarrow I\)}
        \UIC{\(\vdash \la x. \la y. (x, y) : A \rightarrow B \rightarrow (A \times B)\)}
    \end{prooftree}
\end{primer}




\pravilo{Vsotni tip}
Vsotni tip ima levo in desno pravilo vpeljave.
%
\begin{equation*}
    \begin{bprooftree}
        \AXC{\(\G \vdash M: A\)}
        \RL{\(+ I_L\)}
        \UIC{\(\G \vdash \inl~M :  A + B\)}
    \end{bprooftree}
    \begin{bprooftree}
        \AXC{\(\G \vdash M : B\)}
        \RL{\(\lor I_D\)}
        \UIC{\(\G \vdash \inr~M : A + B\)}
    \end{bprooftree}
\end{equation*}
%

Ko želimo uporabiti element vsotnega tipa ne vemo, ali priprada levi ali desni strani. Torej moramo znati delati tako z levo, kot tudi z desno stranjo.
\begin{prooftree}
    \AXC{\(\G \vdash M : A + B\)}
%
    \AXC{\(\G, x: A \vdash N_1 : C\)}
%
    \AXC{\(\G, y : B \vdash N_2 :  C\)}
    \RL{\(+ E\) }
    \TIC{\(\G \vdash \case~M~(\inl~x \mapsto N_1 \mid \inr~y \mapsto N_2) : C\)}
\end{prooftree}



\begin{primer}
    Poskusimo pokazati, da so elementi \(A + B\) v korespondenci z elementi \(B + A\) oziroma. Dovolj bo, da pokažemo, da obstaja funkcija tipa \(A + B \rightarrow B + A\).

    
    \begin{prooftree}
        \AXC{?}
        \UIC{\(\vdash ? : (A + B \rightarrow B + A)\)}
    \end{prooftree}
    Potrebujemo pravilo vpeljave za implikacijo
    \begin{prooftree}
        \AXC{?}
        \UIC{\( x: A + B \vdash ? : B + A\)}
        \RL{\(\rightarrow I\)}
        \UIC{\(\vdash \la x.? : (A + B \rightarrow B + A)\)}
    \end{prooftree}
    V kontekstu imamo na voljo \(A + B\). To želimo uporabiti

\begin{center}
    \small
        \begin{prooftree}
            \AXC{}
            \RL{\(Ax\)}
            \UIC{\(x: A + B \vdash x : A + B\)}
            \AXC{?}
            \UIC{\(x: A + B, a: A \vdash ? : B + A\)}
            \AXC{?}
            \UIC{\(x: A + B, b : B \vdash ? : B + A\)}
            \RL{\(+ E\)}
            \TIC{\( x: A + B \vdash \case~x~(inl~a \mapsto ? \mid \inl~b \mapsto ?) : B + A\)}
            \RL{\(\rightarrow I\)}
            \UIC{\(\vdash \la x.\case~x~(\inl~a \mapsto ? \mid \inr~b \mapsto ?) : (A + B \rightarrow B + A)\)}
        \end{prooftree}
\end{center}
    Potrebno je še najti 
    \begin{prooftree}
        \AXC{?}
        \UIC{\(x: A + B, a: A \vdash ? : B + A\)}
    \end{prooftree}
    in simetrično
    \begin{prooftree}
        \AXC{?}
        \UIC{\(x: A + B, b : B \vdash ? : B + A\)}
    \end{prooftree}
    Poskusimo s pravilom vpeljave za vsotni tip.
    \begin{prooftree}
        \AXC{?}
        \UIC{\(x: A + B, a: A \vdash a: A\)}
        \RL{\(+ I_R\)}
        \UIC{\(x: A + B, a: A \vdash \inr~a : B + A\)}
    \end{prooftree}
    Sedaj pa uporabimo še pravilo aksioma.
    \begin{prooftree}
        \AXC{}
        \RL{\(Ax\)}
        \UIC{\(x: A + B, a: A \vdash a: A\)}
        \RL{\(+ I_R\)}
        \UIC{\(x: A + B, a: A \vdash \inr~a : B + A\)}
    \end{prooftree}
    Podobno naredimo tudi za drugo stran. Sedaj lahko vse skupaj združimo v dokazovalno drevo.


    \begin{center}
        \small
        \resizebox{0.8\textwidth}{!}{%
        \vbox{\begin{prooftree}
            \AXC{}
            \RL{\(Ax\)}
            \UIC{\(x: A + B \vdash x : A + B\)}
            \AXC{}
            \RL{\(Ax\)}
            \UIC{\(x: A + B, a: A \vdash a: A\)}
            \RL{\(+ I_R\)}
            \UIC{\(x: A + B, a: A \vdash \inr~a : B + A\)}
            \AXC{}
            \RL{\(Ax\)}
            \UIC{\(x: A + B, b : B \vdash b : B\)}
            \RL{\(+ I_L\)}
            \UIC{\(x: A + B, b : B \vdash \inl~b : B + A\)}
            \RL{\(+ E\)}
            \TIC{\( x: A + B \vdash \case~x~(inl~a \mapsto ? \mid \inl~b \mapsto ?) : B + A\)}
            \RL{\(\rightarrow I\)}
            \UIC{\(\vdash \la x.\case~x~(\inl~a \mapsto ? \mid \inr~b \mapsto ?) : (A + B \rightarrow B + A)\)}
        \end{prooftree}}}
    \end{center}
\end{primer}


